\section{Expose forward and inverse transforms through one entry}
Implement \texttt{h263\_transform8}. Pass an eight-byte-aligned 48-byte request
in \texttt{\$a0} and a sixteen-byte-aligned 128-byte output record in
\texttt{\$a1}. Request fields are unsigned native-endian doublewords:
\begin{center}
  \tagpdfsetup{table/header-rows={1}}
  \begin{tabular}{rll}
    \textbf{Byte offset} & \textbf{Field}  & \textbf{Meaning or range}         \\
    0                    & Input pointer   & Two-byte-aligned signed halfwords \\
    8                    & Input stride    & Even byte count, 16 through 4096  \\
    16                   & Direction       & 0 forward; 1 inverse              \\
    24                   & Mode            & 0 scalar; 1 MSA arithmetic        \\
    32, 40               & Reserved fields & Both zero
  \end{tabular}
\end{center}
The caller supplies eight active halfwords in each of eight mapped rows. The
minimal extent is $7s+16$ bytes. Input need not be sixteen-byte aligned. Output
is disjoint from the request and input. Supply a sixteen-byte-aligned stack and
valid private space, preserving saved general registers and the incoming stack
pointer when returning.

Validate record alignment and non-wrapping extents, input alignment, stride,
direction, mode, reserved fields and the non-wrapping input extent before sample
reads. Invalid metadata returns status one and count zero without reading input
samples or changing output. Then validate all sixty-four signed sample values.
Forward input must lie in $[-256,255]$; inverse input in $[-2048,2047]$.
An out-of-domain sample returns status two and count zero, retaining output.
Sample reads needed for this second validation phase are permitted.

On success, return \texttt{\$v0}=0 and \texttt{\$v1}=64. Write sixty-four signed
native-endian halfwords in eight tight rows of sixteen bytes. Forward order is
$F[v,u]$; inverse order is $f[y,x]$. Preserve all input and request bytes. Stage
results privately until validation and computation complete, so the caller's
record is unchanged after either kind of rejection.

A completed forward result must be within one coefficient unit of the independent
floating-point reference, with the stated clipping. The inverse also requires
one-unit sample accuracy in finite checks, exact zeros for a zero block, and the
separate Annex A accuracy gate described below. Different accurate transforms
need not produce identical samples at every rounding boundary. The public cases
do not substitute for the statistical inverse accuracy test.

Use \path{construction/h263-transform-construction/main.asm}. The provided matrices are
input data for the taught algorithm; a valid alternative may use its own tables
or factorization. Prepare a new working observation without replacing instructions:
\begin{verbatim}
python3 construction/h263-transform-construction/prepare.py \
    --source my-transform.asm --mode 1 --direction 0 \
    --pattern flat --output forward-flat.asm
\end{verbatim}
Direction zero selects forward, one inverse. The flat forward input contains 128
at every sample and its first output is DC 1024. The flat inverse input contains
DC 1024 and no AC terms; its first output is 128. Both completed callers return
count 64 and exit zero. The unfinished entry reports development status four.

The worked forward input is an actual eight-by-eight luma residual from the
previous generated pictures, using current origin $(16,32)$ and reference
$(14,30)$. Its first output coefficient is approximately $-192$.
Worked inverse input combines DC 1024 with several small AC terms; its first
spatial output is approximately 137. A valid implementation may differ by one
at the allowed rounding boundary. The helper installs inputs and typed controls,
retains the complete instruction section and refuses an existing output file.
It does not compute a target function body.

At \texttt{transform\_caller\_return}, before the first caller instruction,
the output is complete and return registers contain zero and 64. The named
status and count fields have not yet received their copies. Inspect the whole
block as signed halfwords; the ordinary printed first value is only a convenient
observation, not the complete accuracy check.

\section{Check inverse accuracy independently}
An inverse can agree closely with its own forward transform while both use the
wrong normalization. Use an independent reference to test the inverse alone.
H.263 Annex A specifies three 10\,000-block input constructions, their sign-reversed
counterparts, a reference forward/inverse calculation and several error measures.
The reference construction uses at least 64-bit floating-point arithmetic; it
also includes source ranges wider than this exercise's normal forward input.
The student forward entry is not used to manufacture the inverse's expected data.

For each output position, measure peak error, mean squared error and signed mean
error across a cohort. Also measure overall mean squared and signed mean errors.
A zero coefficient block must reconstruct exactly zero. The prescribed generator's
32-bit state and its reset/continuation convention must be recorded with the data.
The retained check uses initial seed one, continuous state across the three base
ranges and the same saved samples negated for each counterpart.

The actual target runs covered both implementations, both modes and both byte
orders: 480\,000 inverse blocks and 48 additional zero calls. Every cohort passed.
The table gives the worst observed values across those runs for each precision:
\begin{center}
  \tagpdfsetup{table/header-rows={1}}
  \begin{tabular}{lrrr}
    \textbf{Measure}            & \textbf{Limit} & \textbf{Scale $2^{16}$} & \textbf{Scale $2^{15}$} \\
    Peak magnitude              & 1              & 1                       & 1                       \\
    MSE per position            & 0.06           & 0.0028                  & 0.0037                  \\
    Overall MSE                 & 0.02           & 0.0016640625            & 0.0026171875            \\
    Mean magnitude per position & 0.015          & 0.0012                  & 0.0015                  \\
    Overall mean magnitude      & 0.0015         & 0.000025                & 0.0000390625
  \end{tabular}
\end{center}
Here mean magnitude means the absolute value of the signed mean error, not the
mean absolute error. Keep those statistics distinct. The latter would be a
different measure and cannot replace the prescribed bias checks.

These are inverse-transform accuracy results. They do not establish quantization,
bitstream syntax, final video quality or complete encoder speed. The instructor
accuracy runner retains input hashes, target sources, byte order, mode and all
cohort metrics. The public finite feedback remains useful while constructing the
function, and the full gate is required before using the inverse in reconstruction.
